\documentclass[11pt]{article}

\usepackage[T1]{fontenc}
\usepackage{lmodern}
\usepackage[margin=1in]{geometry}
\usepackage{amsmath,amssymb,amsthm}
\usepackage{booktabs}
\usepackage{tabularx}
\usepackage{array}
\usepackage{microtype}
\usepackage{xcolor}
\usepackage{enumitem}
\usepackage{hyperref}

\hypersetup{
  colorlinks=true,
  linkcolor=blue!55!black,
  citecolor=blue!55!black,
  urlcolor=blue!55!black,
  pdftitle={Exact d=2 commutativity in a finite strong-operator CPOBC system},
  pdfauthor={To-Marigi}
}

\setlist{nosep}
\setlength{\parindent}{0pt}
\setlength{\parskip}{0.45em}
\newtheorem{theorem}{Theorem}
\newtheorem{proposition}{Proposition}
\newtheorem{lemma}{Lemma}
\newcommand{\QQ}{\mathbb{Q}}
\newcommand{\KK}{\mathbb{K}}
\newcommand{\GL}{\mathrm{GL}}
\newcommand{\Cpobc}{\mathrm{CPOBC}}
\newcommand{\artifact}[1]{\texttt{\small #1}}

\title{\textbf{Exact \(d=2\) commutativity in a finite\\
strong-operator CPOBC system}}
\author{To-Marigi\\
\small Independent researcher}
\date{Final-Theory Bench v0.3.7 -- 30 July 2026}

\begin{document}
\maketitle

\begin{abstract}
We give a computer-assisted theorem for a finite, explicitly defined
strong-operator model of causal-past-ordered Bell causality (CPOBC).
The model is the repository's nonsingular \(n\leq 4\) compiler, not the
unrestricted quantum sequential-growth algebra.  In dimension two, every
solution of its literal, printed-\(Q_{n+1}\) branch has
\([Q_i,Q_j]=0\) for \(1\leq i<j\leq 4\).  The proof removes \(Q_5\)
before elimination: 2,564 canonical numerators split exactly into 2,552
\(Q_5\)-free and 12 \(Q_5\)-dependent expressions, and every literal
solution satisfies the weaker 2,552-expression shared core.  A complete
three-stratum cover of the commuting ratios \(Q_1^{-1}Q_i\), \(i=2,3,4\),
reduces the noncommutative alternative to 21 exact rational saturation
calculations.  All 21 close over \(\QQ\).  Two 21-chart prime-field
campaigns agree but are used only as scouts.

As a separate proposition, we verify a scalar-chain open family in which
\(Q_2,Q_3,Q_4\) are scalar multiples of \(Q_1\) and the literal \(Q_5\)
fiber is \(\GL_2\).  That proposition is established by exact evaluation
of 1,001 direct reduced-operator relations plus exhaustive, hash-bound
compiler-provenance coverage of all 2,564 canonical numerators.  It is not
a second independent implementation and it does not classify \(Q_5\) over
an arbitrary base.  Machine-readable certificates, budgets, and negative
claim boundaries are included.
\end{abstract}

\section{Problem and claim boundary}

Classical sequential growth combines discrete general covariance with a
Bell-causality condition on causal-set births
\cite{RideoutSorkin2000}.  Srivastava and Surya formulate operator-valued
quantum Bell-causality orderings and study the transition algebra generated
by nonsingular operators \(Q_n\) \cite{SrivastavaSurya2026}.  Their
time-ordered and non-time-ordered versions collapse to commutative
transition algebras.  Their CPOBC ordering instead yields additional
relations and leaves the general representation problem open.  The
displayed \(d=2\) computation in their Appendix starts from a
Pauli-proportional ansatz, Eq.\ (165), and proves that this particular
ansatz does not extend to the full algebra.  It is not a classification of
arbitrary invertible \(2\) by \(2\) matrices.

There is also a source-index ambiguity that must remain visible.  On PDF
page 27, atomisation Eq.\ (112) is followed by printed Eq.\ (113) with
\(Q_{n+1}\).  The Appendix's explicit path calculation, Eq.\ (163), uses
the stage-\(n\) generator instead.  The repository therefore keeps two
source-index branches:
\begin{enumerate}[label=(\roman*)]
  \item the \emph{literal branch}, which materialises printed \(Q_{n+1}\);
  \item the \emph{derived branch}, which materialises the \(Q_n\) index
        obtained from the preceding atomisation formula.
\end{enumerate}
No assertion about author intent is made.

The second semantic boundary concerns operator identities.  The source
paper introduces general covariance on a distinguished initial vector and
uses reachable-state reasoning before later atomisation steps employ
operator equalities.  The repository's
\artifact{PAPER\_STRONG\_OPERATOR\_PROFILE} materialises the selected GC,
MSR, CPOBC, and path relations as strong operator identities.  This is an
explicit project formalisation and strengthening; it is not described here
as a uniform hypothesis of Ref.\ \cite{SrivastavaSurya2026}.

Our main statement is consequently finite and conditional on that
formalisation.

\begin{theorem}[Finite literal-branch \(d=2\) commutativity]
\label{thm:main}
Let \(\KK\) be an algebraically closed field of characteristic zero.  Let
\[
  (Q_1,Q_2,Q_3,Q_4,Q_5)\in\GL_2(\KK)^5
\]
satisfy the literal branch of the frozen \(n\leq4\), nonsingular,
\artifact{PAPER\_STRONG\_OPERATOR\_PROFILE} system released by
Final-Theory Bench v0.3.7.  Then
\[
  [Q_i,Q_j]=0 \qquad (1\leq i<j\leq4).
\]
\end{theorem}

The conclusion intentionally excludes \(Q_5\), the transition algebra
outside the compiled inventory, weak or occurrence-dependent operator
semantics, \(d\geq3\), and the infinite-stage CPOBC algebra.

\section{Frozen finite system and the
  \texorpdfstring{\(Q_5\)}{Q5}-free implication}

The compiler freezes reduced operator relations and scalarises their
\(2\) by \(2\) matrix entries after replacing every inverse by its
adjugate divided by its determinant.  This is a calculation in a
localised commutative coordinate ring; adjugate identities and
localisation are standard \cite{StacksAlgebra,MaraisRen2015}, while exact
saturation is a standard ideal operation \cite{BerthomieuEtAl2023}.
These general references do not certify the project-specific relation
inventory or any chart result.

For each source-index branch the frozen direct system contains 1,001
matrix relations.  The literal scalar system contains 2,564 distinct
canonical numerator records.  Version 0.3.6 and the v0.3.7 partition
certificate give the exact disjoint split
\[
  2564 = 2552\ \text{\(Q_5\)-independent}
       + 12\ \text{\(Q_5\)-dependent}.
\]
All 191 frozen denominator records are \(Q_5\)-free.  At the matrix-relation
level, the \(Q_5\)-free subsystem is the 976-relation shared core
\[
  976 = 700\ \Cpobc + 255\ \mathrm{GC} + 21\ \mathrm{MSR}.
\]
The 25 literal path-consistency relations are kept outside that core.

\begin{lemma}[Forward implication]
\label{lem:forward}
Every solution of the 2,564-numerator literal system satisfies the
2,552-numerator shared core.  Therefore, if the shared core forces
\([Q_i,Q_j]=0\) for \(i,j\leq4\), the same commutativity statement holds
for every literal solution, independently of \(Q_5\).
\end{lemma}

\begin{proof}
The shared core is an exact subset selected by canonical equation IDs,
expression IDs, provenance families, and the absence of \(Q_5\) from the
symbol inventory.  Restricting a solution to a subset of its equations is
immediate.  Only this direction is used: no equality of solution schemes,
reverse implication, minimal-generator claim, or Groebner-basis claim is
required.
\end{proof}

This small logical point removes the unsupported premise in the older
49-chart campaign.  Version 0.3.5 placed \(R_5=Q_1^{-1}Q_5\) in a
simultaneous standard form with \(R_2,R_3,R_4\), although the compiled
\(n\leq4\) inventory did not force \([R_5,R_j]=0\).  Those frozen 49
charts and 225 backend runs remain valid inside that ansatz, but they are
not used as a cover of all literal solutions.  The v0.3.7 proof never
introduces \(Q_5\) or \(R_5\).

\section{Complete ratio strata in dimension two}

Set
\[
  R_i=Q_1^{-1}Q_i,\qquad i=2,3,4.
\]
The all-stage source relation Eq.\ (120), specialised only to indices
within \(1,\ldots,4\), has the form
\[
  Q_nQ_1^{-1}Q_m=Q_mQ_1^{-1}Q_n .
\]
Because \(Q_1\) is invertible, this is equivalent to
\([R_n,R_m]=0\).  The relevant instances occur in the frozen shared core.
No index-5 instance of the source theorem is imported.

Over an algebraically closed characteristic-zero field, a commuting triple
of \(2\) by \(2\) matrices has the following exhaustive form, up to
simultaneous similarity:
\begin{description}
  \item[S1: distinct-spectrum pivot.] Some \(R_i\) has two distinct
  eigenvalues.  It is diagonal, and every matrix commuting with it is
  diagonal.
  \item[S2: common nilpotent direction.] There is no distinct-spectrum
  member, but some \(R_i\) is nonscalar.  Every member is
  \(\lambda I+\mu N\) for a common nonzero \(N\) with \(N^2=0\).
  \item[S3: scalar triple.] Every \(R_i=\lambda_i I\).
\end{description}
The general simultaneous-similarity theory of \(2\) by \(2\) tuples is
well known \cite{Florentino2009,Florentino2006}.  The particular
S1/S2/S3 normal forms, localisation factors, and finite chart split used
here are project derivations and are tested against the frozen compiler.

Choosing the first eligible pivot and covering the residual zero/nonzero
conditions produces 12 S1 charts and 9 S2 charts.  S3 needs no elimination:
\[
  Q_i=Q_1R_i=\lambda_iQ_1\quad(i=2,3,4),
\]
so \(Q_1,\ldots,Q_4\) commute structurally.  The S1/S2 charts encode
noncommutativity by saturating each independent component of the six
commutators \([Q_i,Q_j]\), \(1\leq i<j\leq4\).

\section{Exact chart calculation}

For every S1/S2 chart, the 976 direct reduced-operator relations are
evaluated over a rational function field and converted to a polynomial
ideal over \(\QQ\).  The recorded determinant and chart factors define
the open locus.  Sequential Sage/Singular saturation then asks whether a
nonzero commutator component can survive.  A unit ideal is an exact
certificate that the corresponding localised branch is empty.  Since a
unit ideal over \(\QQ\) remains a unit ideal after extension of scalars,
the result applies to the characteristic-zero field \(\KK\) in
Theorem~\ref{thm:main}.

\begin{table}[ht]
\centering
\caption{Proof campaign.  Prime fields are scouts, not proof fields.}
\label{tab:campaign}
\begin{tabular}{lrrrr}
\toprule
Field & S1 charts & S2 charts & Completed & Role\\
\midrule
\(\QQ\)       & 12 & 9 & 21 & exact proof\\
\(\mathrm{GF}(32003)\) & 12 & 9 & 21 & scout\\
\(\mathrm{GF}(32009)\) & 12 & 9 & 21 & scout\\
\midrule
Total & 36 & 27 & 63 & \\
\bottomrule
\end{tabular}
\end{table}

All 21 rational runs have
\artifact{exit\_status=COMPLETED},
\artifact{proof\_eligible=true}, complete localisation, complete
commutator-component coverage, and memory use within the externally owned
8 GiB limit.  In S1, nine charts are empty and three have no surviving
noncommutative component.  In S2, three charts are empty and six have no
surviving noncommutative component.  The two scout fields reproduce all
21 chart verdicts.  They are not promoted to characteristic-zero
certificates.

\begin{proof}[Proof of Theorem~\ref{thm:main}]
Let a literal solution be given.  By Lemma~\ref{lem:forward}, its
\(Q_1,\ldots,Q_4\) coordinates satisfy the \(Q_5\)-free shared core.
The ratio relations make \(R_2,R_3,R_4\) a commuting triple.  The
three-stratum lemma places it in S1, S2, or S3 after simultaneous
similarity.  Every S1 and S2 noncommutative branch has an exact
characteristic-zero unit-ideal certificate; S3 is structurally
commutative.  Hence every commutator among \(Q_1,\ldots,Q_4\) vanishes.
\end{proof}

\section{A scalar-chain \texorpdfstring{\(Q_5\)}{Q5} fiber}

The main theorem says nothing about \(Q_5\).  A separate exact identity
calculation shows why that omission matters.

\begin{proposition}[Literal scalar-chain open family]
\label{prop:fiber}
Let
\[
\begin{gathered}
 Q_1=\begin{pmatrix}a&b\\c&d\end{pmatrix},\qquad
 Q_5=\begin{pmatrix}e&f\\g&h\end{pmatrix},\\
 Q_i=\lambda_iQ_1\qquad(i=2,3,4).
\end{gathered}
\]
On the certificate-defined open set where the 15 recorded irreducible
factors are nonzero, all 1,001 frozen literal reduced-operator relations
vanish identically.  Fourteen factors involve only the base coordinates
\(a,b,c,d\) and \(\lambda_2,\lambda_3,\lambda_4\).  The only
\(Q_5\)-dependent factor is \(fg-eh\).  Thus, above every base point in
the corresponding 14-factor open set, the literal \(Q_5\) fiber is
\(\GL_2\).
\end{proposition}

The proof has two distinct components, but only one actual evaluation
route.  First, Sage 10.9 evaluates all 1,001 direct reduced-operator
matrix relations over the exact \(\QQ\) fraction field.  In both the
literal and derived branches, all 1,001 relation IDs are selected and
evaluated, all residual entries are zero, and all 165 reconstructed
transition predicates are checked.  These runs carry
\artifact{identity\_proof\_eligible=true}; the separate ideal-solve flag
\artifact{proof\_eligible} is correctly false.

Second, a static compiler certificate proves that the direct inventory
covers the scalar inventory.  For each branch:
\[
  1001 = 979\ \text{represented relations}
       + 22\ \text{identically-zero path relations}.
\]
The 979 represented relations contribute all four matrix entries, giving
3,916 provenance records that cover all 2,564 canonical numerators.
The represented relation IDs are exactly the complement of the 22 path
IDs.  Every source stage, family, branch, original-expression ID, and
v0.3.3 residual digest is cross-checked.  The compact arena is schema- and
SHA-256-bound to the frozen polynomial system, and every referenced target
index is in range.

This establishes the proposition without claiming an actual second
compact-arena evaluation of 2,564 expressions.  The two components share
compiler and backend lineage and are not independent implementations.
The proposition constructs a noncommutative open family whenever \(Q_5\)
is chosen not to commute with \(Q_1\), but only over the scalar-chain base.
It neither classifies \(Q_5\) above an arbitrary literal solution nor
produces a representation of the intended infinite CPOBC algebra.

\section{Reproducibility and integrity}

Table~\ref{tab:artifacts} gives the proof-decisive files and selectors.
Every backend response is bound to the external budget SHA-256, the full
semantic request digest, the selected relation/equation digests, the
coefficient field, and the memory limit.  The v0.3.7 scope addendum
re-opens all 63 Phase-1 certificates, recomputes their hashes and proof
gates, reconstructs the two scout maps, and refuses a strong verdict if
any chart is absent or ineligible.

\begin{table}[ht]
\centering
\caption{Machine-readable claim map.}
\label{tab:artifacts}
\small
\begin{tabularx}{\textwidth}{@{}
  >{\raggedright\arraybackslash}p{0.30\textwidth}
  >{\raggedright\arraybackslash}X@{}}
\toprule
Artifact & Decisive fields\\
\midrule
\path{results/v0.3.7_q5_free_partition.json}
& \(2564=2552+12\), selected equation/relation digests, no \(Q_5\) in
  the shared core\\
\path{results/v0.3.7_q5_free_elimination.json}
& 21 exact QQ runs, 42 scouts, chart coverage, verdict
  \artifact{LITERAL\_Q1\_Q4\_COMMUTATIVITY\_PROVED}\\
\path{results/v0.3.7_scope_addendum.json}
& frozen v0.3.4--v0.3.6 hashes, 63 certificate bindings, corrected
  historical claim scope\\
\path{results/v0.3.7_general_scalar_chain_fiber.json}
& 1,001 exact identities per branch, 2,564-numerator provenance
  coverage, 15/14 factor domains, identity proof gate\\
\path{config/v0.3.7_budget.json}
& 3,600 s per chart, 43,200 s total, 8 GiB; no runtime defaults\\
\bottomrule
\end{tabularx}
\end{table}

The budget file SHA-256 is
\artifact{f192a37b97e70c84aeae594db75832bdcf632bf00a78009ff45f09b1bcca2e50}.
The exact calculation used Sage 10.9 and Singular 4.4.1.  The scalar-chain
aggregate points only to the final request-bound certificates; superseded
and timed-out exploratory certificates are excluded from the release
inventory.

\section{Limitations and conclusion}

The result resolves a precise finite question: within the repository's
nonsingular \(n\leq4\) strong-operator literal system, \(d=2\) forces
pairwise commutativity of \(Q_1,\ldots,Q_4\).  It does not settle:
\begin{itemize}
  \item the intended index in printed Eq.\ (113);
  \item weak, fixed-vector, or reachable-state operator semantics;
  \item a restriction or lifting theorem between the frozen compiler and
        the full finite- or infinite-stage CPOBC algebra;
  \item the separately identified general Eq.\ (139) materialisation gap;
  \item \(Q_5\) over a general base, dimensions \(d\geq3\), or singular
        transitions;
  \item extension of an operator or vector measure to a completed event
        sigma algebra.
\end{itemize}

Accordingly, the repository-wide scientific verdict remains
\artifact{FINAL\_THEORY\_OPEN}.  The finite theorem is stronger and cleaner
than the old simultaneous-\(R_5\) ansatz because it uses only a necessary
\(Q_5\)-free subsystem.  The scalar-chain proposition simultaneously shows
that the literal source-index branch can have a free, generically
noncommuting \(Q_5\) fiber on a special open base.  Keeping these two
statements separate is essential: the first is a no-go theorem for
\(Q_1,\ldots,Q_4\) in the frozen literal system, while the second is a
bounded construction that exposes the remaining source-index sensitivity.

\section*{AI-use disclosure}

The author used OpenAI Codex to assist with code review, exact-computation
orchestration, test generation, bibliography organisation, and language
editing.  The human author selected the research scope and budget and is
responsible for every claim.  AI systems are not listed as authors.  All
proof-bearing claims are tied to exact, inspectable artifacts and tests.

\section*{Data and code availability}

Source code, frozen inputs, exact certificates, and the paper source are
contained in the \emph{universe-theory-lab} repository on the
\artifact{codex/final-theory-v0.3.7-publication-20260729} branch.  No DOI
is claimed in this draft.  Release and archival metadata are supplied
separately and must not be read as an external publication.

\bibliographystyle{plain}
\bibliography{references}

\end{document}
